\section{Discussion and Limitations}
\label{sec:discussion}

\paragraph{What the proof-of-concept demonstrates.}
The PoC results in Table~\ref{tab:preliminary} should be read as a mechanism validation, not a deployment claim. The 89.3\% precision achieved by EVICT (No Symb.) is conditional on a specific experimental scope: three CWEs (CWE-89, CWE-78, CWE-190) chosen because their preconditions are well-localized to the extracted slice, a single analyzer (SpotBugs), and a synthetic benchmark (Juliet) with clean ground-truth labels derived from manifests. Within this scope, the ablation structure tells a coherent story: evidence conditioning improves precision by 5.2 points and reduces ECE at full coverage; conformal calibration further reduces selective risk by 35\% (0.165\,$\to$\,0.107) by rejecting uncertain cases; and symbolic escalation safely recovers 2.3 points of coverage (87.3\%\,$\to$\,89.6\%) while adding 1.9 points of precision. Each component earns its operational overhead. The open question---addressed in the planned full evaluation---is whether these gains transfer to harder CWEs, cross-project settings, and real-world alert distributions.

\paragraph{Interpreting the multi-model baseline.}
The multi-model zero-shot study (Table~\ref{tab:multimodel}) reveals a consistent precision--recall trade-off across the lite-model landscape. Claude Haiku 4.5 occupies the conservative end: it aggressively classifies most borderline alerts as FP, achieving 80.3\% FP mitigation but retaining only 36.8\% of true positives. GPT-4o Mini occupies the opposite extreme: near-universal acceptance (98.8\% coverage) with minimal discrimination (13.0\% FP mitigation). Neither profile is suitable for production use without a calibration layer. The key insight is that EVICT's conformal selective prediction does not require a model that is accurate on every alert; it requires a model whose confidence scores are informative about which alerts it can and cannot correctly resolve. The PoC results suggest this condition holds for taint-style CWEs even for lite models: calibrated confidence correctly identifies which alerts are reliably classifiable. What the multi-model study also shows is that the conformal guarantee cannot recover semantic understanding that the model lacks---for CWEs requiring knowledge of value ranges (CWE-190), cross-method aliasing, or library-specific sinks, even well-calibrated abstention cannot substitute for stronger reasoning. This is the motivation for advanced prompting (chain-of-thought, few-shot) and the full evaluation.

\paragraph{Limitations and scope.}
The primary limitation is empirical scope. The current study uses a single analyzer, a synthetic benchmark, and only three CWE classes, none of which involve concurrency, reflection-heavy code, or cross-domain semantics. Conformal validity holds under exchangeability between calibration and test data, an assumption that is weakened by cross-project shift, analyzer version changes, and evolving codebases. The asymmetric label-conditional extension (Theorems~\ref{thm:label_conditional} and~\ref{thm:exact_opt}) has not yet been validated empirically; the current implementation instantiates only the symmetric split-conformal special case. No head-to-head comparison against IRIS or the strongest recent baselines has been conducted. A further limitation is that Table~\ref{tab:multimodel} (multi-model baseline) reports single-run aggregate metrics without fold-level cross-validation; variance estimates for that table are deferred to the full evaluation.

Manual inspection of 50 residual errors in the accepted set reveals four dominant failure modes: sanitization guards located outside the 50-line slice boundary (the most common), aliased control flow that is not captured by the backward slice, reflection-based dispatch that prevents slice completion, and Z3 path constraints that exceed solver capacity within the 2s timeout. Each failure mode corresponds to a case where the EvidencePack is genuinely incomplete---suggesting that tighter slice-completeness checks and adaptive slice bounds could further reduce precision losses on accepted alerts.

Sensitivity of the schema-guided prompt to minor perturbations (reworded preconditions, alternative step ordering) remains an open ablation. The planned evaluation includes a prompt-robustness sweep to quantify this sensitivity and determine whether few-shot exemplars narrow the variance across models. Our overall position is therefore narrower than a deployment claim: EVICT demonstrates a coherent and theoretically grounded decision pipeline, with preliminary evidence that the mechanism works as intended under controlled benchmark conditions. The expanded evaluation in Section~\ref{sec:evaluation} is necessary to establish whether these gains transfer to real-world complexity.
